\section{Introduction}\label{sec:introduction}

A \emph{Kakeya set} in \(\R^n\) is a set containing a unit line segment
in every direction. The Hausdorff-dimension Kakeya conjecture asserts
that every such set has Hausdorff dimension \(n\). We prove its
four-dimensional case.

\begin{theorem}\label{thm:main}
Let \(K\subset\R^4\). Suppose that for every \(e\in S^3\) there is a point
\(a_e\in\R^4\) such that
\[
  \{a_e+te:0\leq t\leq1\}\subset K.
\]
Then \(\dim_H K=4\).
\end{theorem}

Theorem~\ref{thm:main} resolves the Hausdorff-dimension Kakeya conjecture
positively in four dimensions. Neither \(K\) nor the choice of its
witnessing segments is required to be measurable. In particular, no
packing-dimension or stickiness hypothesis is imposed on the family
of witnessing lines.

For arbitrary sets, we use the covering definition
\[
 \mathcal H_\delta^s(K)
 =\inf\left\{\sum_i(\diam U_i)^s:
        K\subset\bigcup_i U_i,\quad \diam U_i\leq\delta\right\},
 \qquad
 \mathcal H^s(K)=\lim_{\delta\downarrow0}\mathcal H_\delta^s(K),
\]
where the covers are countable, and
\(\dim_H K=\inf\{s:\mathcal H^s(K)=0\}\).
This definition requires no measurability assumption. The reduction
in Section~\ref{sec:framework} uses finite selections of witnessing
segments and outer-measure estimates in direction space.

Related dimension statements and geometric consequences are developed
in Section~\ref{sec:consequences}. Besides the packing- and
bounded-set Minkowski-dimension implications, these include a lower
bound in higher dimensions by orthogonal projection, arbitrary
direction sets and extension of segments to lines, and Nikodym and
curved-Kakeya applications. The latter implications use the transfer
results of \citet{KeletiMatheEquivalences},
\citet{GaoLiuXi2025}, and \citet{NadjimzadahLifting}, with their
distinct hypotheses stated there.

\subsection{Context and input estimates}

The Kakeya problem connects the geometry of line segments with
estimates for oscillatory operators. The classical needle problem
asks how small a planar region can be if a unit segment is to turn
continuously within it and return with its endpoints reversed; the
area can be arbitrarily small \citep{Besicovitch1928,Davies1971}.
A central difficulty in the dimension problem is that many segments
can overlap substantially even when their directions are well
distributed. Besicovitch's constructions show that a Kakeya set can
have Lebesgue measure zero, whereas Davies proved full Hausdorff
dimension in the plane \citep{Besicovitch1928,Davies1971}.

For bounded Kakeya sets, the upper Minkowski formulation asks for
dimension \(n\), measured by covering numbers at one common scale;
the Hausdorff formulation allows covers with varying diameters.
For fixed ambient dimension \(n\ge2\), let \(T_\delta(a,e)\) be the
cylinder centered at \(a\in\R^n\), with axis direction
\(e\in S^{n-1}\), length one, and transverse radius \(\delta\).
Define the maximal operator by
\(K_\delta f(e)=\sup_{a\in\R^n}|T_\delta(a,e)|^{-1}
\int_{T_\delta(a,e)}|f(x)|\,dx\).
The stronger Kakeya maximal conjecture asks that, for every
\(\varepsilon>0\), there be a constant \(C_{\varepsilon,n}\) such that
\[
 \|K_\delta f\|_{L^n(S^{n-1})}
 \leq C_{\varepsilon,n}\delta^{-\varepsilon}
       \|f\|_{L^n(\R^n)}
\]
for all \(f\in L^n(\R^n)\) and \(0<\delta<1\)
\citep{KatzTaoNewBounds}. For partial maximal estimates, the
dimension parameter records the Hausdorff lower bound supplied by
that estimate. Theorem~\ref{thm:main} proves the four-dimensional
Hausdorff assertion.

Wolff's bound \(\dim_H K\ge(n+2)/2\) gives dimension at least three
in \(\R^4\). His geometric method studies the hairbrush of a tube:
the collection of tubes meeting it
\citep{Wolff1995,GuthZahlPolynomialWolff}.
Bourgain connected line incidences in separated slices to sumset
and difference-set estimates, using Gowers's quantitative form of
the Balog--Szemer\'edi theorem \citep{Bourgain1999}.
Katz and Tao developed sums--differences iterations for
higher-dimensional bounds \citep{KatzTaoNewBounds}.
In three dimensions, Katz, \L{}aba, and Tao combined multiscale
clustering and local planar structure with the arithmetic method
to improve the upper Minkowski lower bound beyond \(5/2\)
\citep{KatzLabaTao2000}.
\L{}aba and Tao then used multiscale geometric structure to obtain
upper Minkowski dimension strictly greater than \((n+2)/2\) for
\(n\ge4\), in particular greater than three in four dimensions
\citep[Theorem~1.3]{LabaTaoMediumDimension}.

The later four-dimensional Hausdorff improvements addressed
concentration near algebraic sets. For compact Kakeya sets,
the polynomial Wolff estimates of Guth and Zahl, together with the
algebraic direction bounds of Zahl and Katz and Rogers, give the
lower bound \(3+1/40\)
\citep{GuthZahlPolynomialWolff,ZahlSeveri,KatzRogersPolynomialWolff}.
The polynomial Wolff axioms restrict concentration near low-degree
algebraic sets; the accompanying multilinear estimates exploit
quantitative linear independence of tube directions
\citep{GuthZahlPolynomialWolff}.
Katz and Zahl's planebrush organizes tubes around a common plane
and combines this geometry with Guth and Zahl's trilinear estimate.
It gives Hausdorff dimension at least \(3.059\), while their distinct
maximal-function estimate has dimension parameter at least \(3.049\)
\citep{KatzZahlPlanebrush}.
Borges, Chan, Chen, Liu, Xi, and Zhan subsequently improved the
maximal parameter to
\((159+\sqrt{145})/56\approx3.0543\), combining the planebrush
with a planar two-ends Furstenberg estimate
\citep[Theorem~1.2]{BorgesEtAlRestrictionKakeya}.
This improves the maximal estimate while leaving the larger
Hausdorff bound \(3.059\) as a separate conclusion.

Rai Choudhuri proved the bound \(13/4\) for compact sticky sets,
which admit a witnessing line family of packing dimension three
\citep{RaiChoudhuriSticky}. Wang and Zakharov have since obtained a
full-dimensional union estimate for almost AD-regular sticky tube
families satisfying convex Wolff axioms and a dense-shading condition
\citep{WangZakharovSticky}; those hypotheses do not cover arbitrary
Kakeya sets. In three dimensions, Wang and Zahl proved the Kakeya
set conjecture \citep{WangZahlThreeDimensional}, and Guth, Wang,
and Zahl subsequently streamlined the proof
\citep{GuthWangZahlStreamlined}.

The three-dimensional estimate used in our reduction is the weighted
full-time plank bound of \citet[Lemma~2.3]{ThreeDimensional}, recorded
as \Cref{lem:weighted-planks}. It is proved there by weighted convex
clustering from the union estimate of
\citet[Theorem~1.1]{GuthWangZahlStreamlined}. It supplies spatial
boxes with enough incidence mass to support the quadratic local
fits described below.
The further analytic inputs have distinct roles. Determinant
multilinear Kakeya \citep{CarberyValdimarsson} controls configurations
of independent velocities. The restricted-triple dot-product
theorem of \citet{WangZahlSticky} and the multiplicative-convolution
estimate of \citet{OrponenDeSaxceShmerkin} provide planar rigidity.
The latter is applied in \Cref{re:retained-edge-rigidity} to construct
scalar coefficients with small interval probabilities.
That step also uses a graph Balog--Szemer\'edi--Gowers argument and
fixed-iterate sumset calculus
\citep{SudakovSzemerediVu2005,TaoVu2006}. The graph argument retains
a substantial set of original incidence edges, which are needed
to fit the resulting planar structure back to the trajectories.

\subsection{Mechanisms of the proof}

The proof uses quadratic local fits to capture line concentration
that persists across scales. Its basic object is a \emph{chart}:
a collection of line pieces in a moving spatial box and a time
interval, on which one quadratic polynomial test stays small.
Each chart also carries a mass requirement, measured using selected
time-line pairs. This prevents a fit to an exceptional collection of
negligible weight from counting as progress. Normalizing inside a
chart produces another problem of the same kind, with its time
interval rescaled to unit length. Charts can therefore be nested.

The two quantities to balance are the thickness of the quadratic
relation and the length of the line pieces on which it holds.
Longer intervals at a prescribed thickness give better charts.
The reduction in Section~\ref{sec:framework} turns a hypothetical
Hausdorff deficit into a weighted line problem with two competing
properties: refining to terminal depth loses less than one power of
density per unit depth, yet every substantial terminal chart system must use
intervals of a definite shrinking length. The proof rules out this
combination.

More precisely, write the lines as
\((t,y_0+tV,w_0+tW)\), where \(y\in\R^2\).
An observation retains exact \((t,y)\) and bins the remaining
coordinate \(w\); increasing the depth refines this bin.
At resolution \(N_0\to\infty\), the selected incidence density in a
bin has a common typical size \(n(r)=N_0^{-f(r)}\).
Here density is measured per \(dt\,dy\) and per hidden bin within
each component of the weighted mixture. The counterexample gives
conditional velocity nonconcentration and
\[
 f(L)-f(r)\le d(L-r)\qquad(0\le r\le L),\qquad d<1,
\]
where \(L\) is the terminal depth.
A chart is called \emph{true} when its incidence mass, divided by
its interval length, is at least this density times its spatial
volume, up to subpower factors.
A chart system's coverage is the selected incidence mass, and mass
\(N_0^{-o(1)}\) is called substantial.
Writing a common interval length as \(N_0^{-h+o(1)}\), its
\emph{time cost} is \(h\).
The \emph{horizon} \(H(E)\) is the infimum of these costs over true
terminal systems of substantial coverage in the weighted problem
\(E\), allowing further selections and subsequences.
The reduction yields \(H(E)>0\).

The exponent below one makes a comparison principle available.
Local graph representations of the polynomial tests give scalar
sheets, which are analytic branches of polynomial equations.
The density bound and the lower mass bounds of retained entries
supply lists of size at most \(\delta^{-d+o(1)}\) at accuracy
\(\delta\), for a fixed \(d<1\).
Under the derivative and sampling bounds in
\Cref{lem:two-color-matching}, two such lists cannot frequently
agree in value while keeping their rescaled derivative data apart.
Alternating fresh conditional samples of sheets and times turns a
persistent derivative discrepancy into a polynomial endpoint map
whose image has more volume than the density bounds allow.
Successful paths remain unnormalized submeasures, so this volume
argument retains their dependence on the preceding choices.
The comparison both lengthens terminal charts and joins the local
approximations constructed later.

To locate a forbidden improvement, we follow nested charts in a
nearly extremal weighted problem. Their relative density and time
profiles become linear, with time rate \(k>0\).
A further extremal choice controls how much thinner the smaller
spatial axis can become than the time scale; the larger width is
comparable to that scale up to subpower factors.
Call the optimized rate of this \emph{narrowness} \(\ell\).
The choices are arranged so that a definite improvement in time,
or an additional gain in narrowness at fixed time when \(\ell\)
is finite, contradicts extremality after completion and return to
the original problem. Section~\ref{sec:critical-paths} gives the
precise order of these choices and the resulting critical paths.

There are three geometric regimes. When \(\ell=0\), enough
separated velocities remain to determine an approximate polynomial
field. Conditional scalar estimates first give polynomial fits
along individual lines. Multilinear Kakeya and interpolation combine
them, and switching between lines makes the field approximately
conservative, producing a potential.
The potential can initially have degree higher than two. The
obstruction has a specific form: the broad directions may lie near
a conic. The proof identifies the surviving cubic and quartic terms
and uses the existing hidden coordinate to encode them, on the
retained trajectories, in a quadratic correction to the chart test.
This restores the degree required for an admissible chart
(\Cref{prop:degree-restoration}).
The same argument with one spatial coordinate also gives zero
horizon for an auxiliary scalar model; that result supplies the
scalar projections used in the other regimes.

When \(0<\ell<\infty\), projection onto the scalar model gives an
alternative. Sufficiently wide projected intervals already yield an
improvement by comparison. Otherwise planar rigidity forces the
trajectories into a horizontal model with equation \(Y'=-Z+tv\).
Its natural boxes have horizontal size \(H\) and transverse size
\(H^2\). The noncommuting horizontal directions force the relevant
scalar sheets to be nearly affine.
At the critical time rate \(2k=1\), a path with three refreshed
line segments fills enough three-dimensional base volume to improve
the fit. Away from that rate the aim is to recover wide projected
intervals. When \(2k<1\), this first requires repeated scalar
estimates that improve the resolution of trajectory parameters over
the whole normalized time interval.

When \(\ell=\infty\), two nested chart labels instead localize
trajectories over the whole normalized time interval.
Normalization can introduce large shifts in the density exponents,
but these cancel in the relative density ratios used by the
argument. A final scalar projection gives charts whose time cost
tends to zero, contradicting the positive cost inherited from the
counterexample.

Each regime initially produces local candidates on substantial
residual portions of the weighted family. To obtain a contradiction,
these fits must become chart assignments with enough total coverage
and with finite lists of labels recoverable from the original
observations. In the finite-narrowness cases, greedy coverage and
the small-list comparison give this passage through
\Cref{prop:candidate-upgrade}. The unbounded-narrowness construction
recovers coverage and finite label lists after its full-time
normalization. In each case, the improved chart system is lifted
to the original weighted problem, where it contradicts the extremal
choice. The comparison, the finite-narrowness conversion, and the
conditional higher-jet argument (\Cref{prop:higher-jets}) are stated
separately because their hypotheses do not require a low-entropy
original line family.

\begin{figure}[t]
 \centering
 \begin{tikzpicture}[
  box/.style={draw,rounded corners=2pt,align=center,inner sep=6pt,
              font=\small,text width=11cm},
  branch/.style={box,text width=4cm,minimum height=15mm},
  flow/.style={-{Latex[length=2mm]},semithick},
  node distance=7mm and 5mm
]
\node[box] (model) {A Kakeya counterexample\\
  \(\longrightarrow\) an exact weighted problem with positive time cost};
\node[box,below=of model] (path) {Critical paths: nested true chart systems\\
  linear density profile, time rate \(k>0\), optimized narrowness \(\ell\)};
\node[branch,below=10mm of path] (finite) {\(0<\ell<\infty\)\\
  wide intervals or a horizontal model\\
  Sections~\ref{sec:finite-narrowness}--\ref{sec:critical-rates}};
\node[branch,left=of finite] (iso) {\(\ell=0\)\\
  polynomial potential;\\
  restoration to degree two\\
  Section~\ref{sec:isotropic}};
\node[branch,right=of finite] (unbounded) {\(\ell=\infty\)\\
  full-time localization and scalar projection\\
  Section~\ref{sec:unbounded-narrowness}};
\node[box,below=11mm of finite] (candidate) {Improvement on every substantial residual\\
  greedy coverage and scalar-sheet matching};
\node[box,below=of candidate] (atlas) {A chart system with recoverable labels in the original problem\\
  smaller terminal time cost or, for finite \(\ell\), excess narrowness};
\draw[flow] (model)--(path);
\draw[flow] (path.south)--(finite.north);
\draw[flow] (path.south) -- ++(0,-4mm) -| (iso.north);
\draw[flow] (path.south) -- ++(0,-4mm) -| (unbounded.north);
\draw[flow] (finite.south)--(candidate.north);
\draw[flow] (iso.south) |- ([yshift=4mm]candidate.north);
\draw[flow] (unbounded.south) |- ([yshift=4mm]candidate.north);
\draw[flow] (candidate)--(atlas);
\end{tikzpicture}

 \caption{The contradiction from a hypothetical Kakeya counterexample.
 A critical weighted path has one of three narrowness regimes.
 Each produces improvements on substantial residuals, which are
 assembled into a chart system with recoverable labels in the
 original problem. In the middle branch, \(2k=1\) uses the three-leg
 switch; \(2k\ne1\) returns to wide intervals, with a bootstrap
 when \(2k<1\). The scalar version of the isotropic argument supplies
 the projection theorem used in the other two branches.}
 \label{fig:proof-map}
\end{figure}

\subsection{Organization and scale conventions}

Section~\ref{sec:framework} gives the weighted models, chart
normalization, and reduction from an arbitrary Kakeya set.
Section~\ref{sec:comparison} proves the scalar-sheet comparison and
terminal extension. Section~\ref{sec:critical-paths} constructs the
critical paths and proves the criterion that turns local candidates
into forbidden improvements.

Section~\ref{sec:scalar-tools} develops the conditional scalar
estimates, and Section~\ref{sec:isotropic} treats the isotropic
case and its auxiliary scalar model.
Section~\ref{sec:retained-edge-tools} proves the planar rigidity
tools and fits their conclusions to the original incidence edges.
Sections~\ref{sec:finite-narrowness} and~\ref{sec:critical-rates}
use these results to treat finite positive narrowness and its
critical time regimes.
Section~\ref{sec:unbounded-narrowness} treats unbounded narrowness
and completes the proof of Theorem~\ref{thm:main}.
Section~\ref{sec:consequences} develops the dimensional and geometric
consequences.

Two scale conventions are essential when these steps are combined.
Section~\ref{sec:critical-paths} distinguishes the original resolution
\(N_0\) from the local tangent resolution \(N\): returned chart
systems must satisfy mass and list bounds subpower in \(N_0\), even
when the local analysis only supplies bounds subpower in \(N\).
Section~\ref{sec:framework} distinguishes current incidence laws
from earlier witness laws. A witness floor used for counting is
paired with a density cap for that same law.
